\documentclass[11pt]{article}
\usepackage[margin=0.85in]{geometry}
\usepackage{lmodern}
\usepackage{microtype}
\usepackage{enumitem}
\usepackage{listings}
\usepackage[T1]{fontenc}

\lstset{
  basicstyle=\ttfamily\small,
  columns=fullflexible,
  keepspaces=true,
  showstringspaces=false,
  breaklines=true,
  frame=single,
  xleftmargin=0.4em,
  xrightmargin=0.4em,
  aboveskip=0.55em,
  belowskip=0.55em
}

\pagestyle{empty}
\setlength{\parskip}{0.35em}
\setlength{\parindent}{0pt}
\setlist[enumerate,1]{leftmargin=1.4em, itemsep=0.75em, topsep=0.3em}
\setlist[enumerate,2]{leftmargin=1.5em, itemsep=0.35em, topsep=0.25em, label=\alph*.}
\setlist[itemize]{leftmargin=1.3em, itemsep=0.15em, topsep=0.2em}

\begin{document}

\begin{center}
{\LARGE\bfseries Midterm 1 --- Sample 1}\\[0.35em]
DATA 412/612
\end{center}

\vspace{0.25em}

Practice exam, questions only. Cumulative through October 14
(weeks 1--7): the Quiz 1 material, plus \texttt{rowwise()} and
\texttt{across()}, \texttt{case\_when()}, \texttt{readr},
exploratory plots, and \texttt{pivot\_longer()}.

Put your answers in an R Markdown document and keep the code visible
(\texttt{knitr::opts\_chunk\$set(echo = TRUE)} in the setup chunk).
You may use your notes, the R help files, and the internet.
Do not discuss the exam with anyone else.
Load the tidyverse. \texttt{mpg} comes with \texttt{ggplot2}.

\begin{enumerate}
\item \textbf{\texttt{case\_when()}} (8 pts).
In \texttt{mpg}, add \texttt{efficiency}: \texttt{"high"} when
\texttt{hwy} is at least 30, \texttt{"mid"} when \texttt{hwy} is at
least 20, and \texttt{"low"} otherwise.
Then, for each \texttt{class} and \texttt{efficiency}, report the
number of cars. Arrange so the largest counts come first.

\item \textbf{Row-wise and \texttt{across()}} (8 pts).
Use this tibble.

\begin{lstlisting}
quizzes <- tibble(
  student = c("Ada", "Lin", "Grace"),
  q1 = c(8, NA, 9),
  q2 = c(7, 10, NA),
  q3 = c(9, 8, 8)
)
\end{lstlisting}

\begin{enumerate}
\item With \texttt{across()}, count the \texttt{NA}s in each quiz column.
\item With \texttt{rowwise()} and \texttt{c\_across()}, add each
student's mean of \texttt{q1}, \texttt{q2}, and \texttt{q3}, ignoring
missing quizzes.
\end{enumerate}

\item \textbf{\texttt{readr} and a plot} (12 pts).

\begin{enumerate}
\item Parse this frame. \texttt{when} should be a date, \texttt{price}
a number (the dollar sign and the word \texttt{n/a} are not numbers),
and \texttt{ok} logical, with the string \verb|NA| meaning missing.
\begin{lstlisting}
raw <- tibble(
  when  = c("2024-03-01", "2024-03-18", "2024-04-02", "2024-04-09"),
  price = c("$12.50", "8", "n/a", "$15"),
  ok    = c("TRUE", "FALSE", "NA", "TRUE")
)
\end{lstlisting}
\item Using \texttt{mpg}, make one plot that compares highway mileage
across \texttt{class}. State which class has the highest median
\texttt{hwy}, and give that median.
\end{enumerate}

\item \textbf{\texttt{pivot\_longer()}} (12 pts).
Reshape \texttt{wide} so there is one row per country and year, a
numeric \texttt{year} column, and a \texttt{cases} column.
Strip the \texttt{y} from the year names. Drop rows whose case count
is missing. Then report the country and year with the most cases.

\begin{lstlisting}
wide <- tibble(
  country = c("A", "B"),
  y2021 = c(10, NA),
  y2022 = c(12, 7),
  y2023 = c(9, 15)
)
\end{lstlisting}
\end{enumerate}

\end{document}
